\documentclass[tikz,border=12pt]{standalone}
\usepackage{amsmath}
\usepackage{kotex}
\usetikzlibrary{arrows.meta,calc}
\definecolor{physical}{RGB}{35,100,175}
\definecolor{stress}{RGB}{198,68,54}
\definecolor{body}{RGB}{225,235,244}
\begin{document}
\begin{tikzpicture}[>=Latex,font=\sffamily,line cap=round,line join=round]
  \node[font=\Large\bfseries] at (7.0,6.2) {실제 물체의 한 점을 응력공간의 한 점으로 나타내기};

  \node[font=\large\bfseries,physical] at (2.1,5.25) {1. 실제 물체};
  \fill[body] (0.3,1.45) rectangle (3.9,4.55);
  \draw[physical,very thick] (0.3,1.45) rectangle (3.9,4.55);
  \draw[physical,thick,->] (0.65,1.8)--(1.45,1.8) node[right] {$\boldsymbol e_1$};
  \draw[physical,thick,->] (0.65,1.8)--(0.65,2.6) node[above] {$\boldsymbol e_2$};
  \draw[stress,thick,fill=white] (2.15,2.8) rectangle (2.75,3.4);
  \fill[stress] (2.45,3.1) circle (3pt);
  \node[stress,above left] at (2.15,3.55) {관심점 $P$};
  \draw[stress,very thick,->] (2.75,3.1)--(3.55,3.1) node[right] {$\sigma_{11}$};
  \draw[stress,very thick,->] (2.15,3.1)--(1.55,3.1);
  \draw[stress,very thick,->] (2.45,3.4)--(3.15,3.4) node[right] {$\sigma_{12}$};
  \draw[stress,very thick,->] (2.45,2.8)--(1.85,2.8);
  \node[align=center,font=\small] at (2.1,0.82) {물체 안의 위치 $P$를 고정하고\\그 점의 응력 상태를 읽는다.};

  \draw[->,very thick] (4.15,3.0)--(5.0,3.0);
  \node[font=\large\bfseries] at (7.0,5.25) {2. $P$점의 응력텐서};
  \node[draw,rounded corners=3pt,very thick,fill=gray!6,minimum width=3.3cm,minimum height=2.5cm,align=center] at (7.0,3.05)
    {$\displaystyle
     [\boldsymbol\sigma]_P=
     \begin{bmatrix}
       80 & 30 & 0\\
       30 & 0 & 0\\
       0 & 0 & 0
     \end{bmatrix}\,\mathrm{MPa}$};

  \draw[->,very thick] (8.95,3.0)--(9.75,3.0);
  \node[font=\large\bfseries,stress] at (11.9,5.25) {3. 응력공간};
  \draw[->,thick] (9.85,1.25)--(14.0,1.25) node[right] {$\sigma_{11}$ (MPa)};
  \draw[->,thick] (9.85,1.25)--(9.85,4.8) node[above] {$\sigma_{12}$ (MPa)};
  \draw[gray!55,dashed] (9.85,3.65)--(13.05,3.65)--(13.05,1.25);
  \fill[stress] (13.05,3.65) circle (3.2pt);
  \node[stress,above right,align=left] at (13.05,3.65) {응력 상태\\$(80,30)$ MPa};
  \node[below] at (13.05,1.25) {$80$};
  \node[left] at (9.85,3.65) {$30$};
  \node[align=center,font=\small] at (11.9,0.55)
    {이 점의 좌표는 물체 안의 위치가 아니라\\점 $P$에서의 응력성분이다.};

  \node[align=center,text=gray!80!black,font=\small] at (7.0,-0.2)
    {하중이 변하면 실제 위치 $P$는 그대로여도 응력공간의 점은 이동하며, 그 이동 궤적이 응력 경로이다.};
\end{tikzpicture}
\end{document}
